\section{线程块内同步、跨线程块同步与可伸缩性}

\subsection{线程块级屏障的语义}

同一线程块内最基本的屏障同步由 \code{__syncthreads()} 提供。其核心语义是：到达该屏障的线程必须等待，直到线程块中所有尚未退出且应参与该屏障的线程都一致地到达同一个屏障实例，然后才能继续执行。

\begin{mechanismbox}
线程块级屏障之所以可以高效实现，是因为同一线程块的全部驻留状态位于同一个流式多处理器中。屏障除了协调控制进度，还会对相关的共享内存和全局内存访问建立必要的顺序关系，使屏障之前的相关写入在屏障之后按照该同步语义被观察。
\end{mechanismbox}

\subsection{发散条件中的屏障为什么危险}

\par\noindent\begin{minipage}{\linewidth}
下面的结构存在根本问题：

\begin{lstlisting}[caption={发散条件中的两个不同屏障实例}]
if (predicate) {
    __syncthreads();
} else {
    __syncthreads();
}
\end{lstlisting}
\end{minipage}\par

表面上看，两条分支都出现了 \code{__syncthreads()}；但如果线程块中的不同线程沿不同分支前进，它们到达的可能是两个不同的静态屏障实例。这样无法保证全部参与线程在同一个屏障上会合，结果可能造成错误甚至死锁。

正确性要求可以概括为：对于某一个线程块级屏障实例，所有应参与的线程必须以兼容、可达且一致的方式到达它。屏障不能依赖会让同一线程块中的线程分成互不相交路径的控制结构。

\begin{warningbox}
“每条分支里都写一个屏障”不等价于“所有线程到达同一个屏障”。判断同步是否正确，必须看线程是否会合到同一个屏障实例，而不是只看源代码中是否都出现了同名函数。
\end{warningbox}

\subsection{为什么线程块必须一次性获得完整驻留资源}

线程块不能只让一部分线程先驻留，剩余线程等以后再分配。原因可以用屏障直接解释：假设线程块前半部分已经驻留并执行到屏障，而后半部分因为资源不足根本无法开始执行，那么前半部分会永久等待后半部分；同时前半部分占据的资源又可能阻止后半部分获得驻留机会，从而形成死锁。

所以，线程块只有在流式多处理器能够为\textbf{整个线程块}提供所需驻留资源时，才会被接纳。资源分配的整体性是线程块级同步语义能够成立的硬件条件之一。

\subsection{普通线程块之间为何不能依赖执行顺序}

如果不同线程块之间没有同步依赖，那么这些线程块可以：

\begin{itemize}
  \item 以任意顺序开始；
  \item 在不同流式多处理器上并行执行；
  \item 在硬件资源较少时部分顺序执行；
  \item 在硬件资源较多时让更多线程块同时执行。
\end{itemize}

这形成了\term{透明可伸缩性}{transparent scalability}：同一份核函数代码可以在具有不同流式多处理器数量的 GPU 上运行，而正确性不依赖具体有多少线程块同时驻留。

\begin{keybox}
线程块执行顺序不受程序员控制，因此普通核函数中的正确性条件必须对所有可能的线程块调度顺序成立。
\end{keybox}

\subsection{跨线程块屏障的死锁风险}

若程序试图让整个网格中的线程块通过普通自制屏障互相等待，问题在于部分线程块可能尚未获得驻留。已经驻留的线程块先到达屏障并占用流式多处理器资源；未驻留的线程块却必须先运行到同一屏障才能解除等待。如果驻留资源被等待中的线程块占满，后续线程块就永远没有机会启动。

受约束的跨线程块同步机制主要包括：

\begin{itemize}
  \item \term{协作组}{cooperative groups}可以在满足特定驻留约束的条件下，支持线程块簇或整个网格范围的屏障；关键前提是相关参与者能够被同时调度到硬件上。
  \item 某些单向生产者—消费者同步可以通过更严格的执行结构建立先后关系，但普通线程块的调度顺序本身仍不能被程序假定。
\end{itemize}

因此，跨线程块同步属于比线程块内同步更高成本、约束更多的机制。除非算法明确需要并满足相应条件，否则应优先保持线程块之间的顺序独立性。
